\documentclass[12pt]{article}

\usepackage[spanish]{babel}
\usepackage[utf8]{inputenc}
\usepackage[T1]{fontenc}
\usepackage{lmodern}
\usepackage[a4paper,margin=2.8cm]{geometry}
\usepackage{setspace}
\usepackage{csquotes}

\setstretch{1.15}

\title{Del premio, del agasajo y de la aceptación}
\author{Vicente Domínguez Arca}
\date{}

\begin{document}

\maketitle

Al menos en el Huniverso que habitamos, esa parte del Universo accesible a los sentidos humanos o a sus extensiones técnicas y tecnológicas, el anacoreta completo no existe. \emph{Anacoreta} es un término con una etimología sugerente: hace referencia a quien se aleja o se retira. Lo sugerente, al menos en mi opinión, descansa en el hecho de que el término haya terminado por designar a aquellas personas que deciden apartarse de la sociedad y llevar una vida solitaria en contacto con la naturaleza que las envuelve.

Algún lector crítico podría decir que unas líneas atrás he realizado una aseveración demasiado rotunda. Me explico. Toda persona, sea anacoreta o no, nace intrínsecamente en el seno de una sociedad; el propio hecho de la existencia se sostiene sobre innumerables aristas que conectan y tejen el tejido social.

Existen entonces dos posibilidades para el anacoreta. O bien nació y se halló ``solo'', o bien creció y maduró en uno o varios grupos sociales antes de decidir alejarse. Al primer caso no puede atribuírsele la definición de anacoreta, pues no se ha alejado porque nunca estuvo cerca: simplemente nació y se halló solo. El segundo caso es lo más próximo que podemos estar del anacoreta completo. Uno puede pretender, e incluso conseguir, alejarse hasta el punto de que los lazos con cualquier otra persona desaparezcan.

Sin embargo, ni siquiera en ese caso extremo puede decirse que alguien sea genuinamente un anacoreta completo. En todo caso podremos afirmar que ha decidido convertirse en anacoreta, puesto que anteriormente estuvo inserto en uno o varios grupos sociales. Así, fuera de toda duda razonable, no parece aceptable sostener que las personas podamos habitar nuestro Huniverso completamente desconectadas de las demás. Somos intrínsecamente seres sociales, y esa parece ser una realidad del Huniverso difícil de rebatir sin recurrir a algún tipo de dogma.

Sugiero ahora al lector desplazarse ligeramente de lo leído hasta aquí, manteniendo, eso sí, su esencia.

Otra etimología que resulta curiosa es la de la palabra \emph{premio}. No creo que sorprenda a nadie su mención; no deja de ser un término completamente convencional. Sin embargo, procede del latín \emph{praemium}. Más allá de la precisión filológica de su origen, el término proyecta una imagen especialmente sugerente: la de aquello que corresponde a quien ha conseguido situarse por delante, a quien, de algún modo, ha llegado antes que los demás. Y es precisamente esa imagen la que me interesa.

El significado actual del término hace referencia a un reconocimiento que se concede a una persona ---no entro aquí a considerar la cuestión de premios concedidos a entornos naturales o a animales--- por entender que ha destacado en algún aspecto. Pero no debemos confundirlo con el premio propio de una competición, que se asemeja más a un trofeo, pues en ese caso las reglas han sido definidas de antemano. El premio al que me refiero responde más bien a algo que, \emph{a posteriori}, se adjudica a una persona porque se considera que ha sobresalido respecto de un criterio cuya relevancia queda plenamente establecida en el propio acto de otorgarlo.

La pregunta resulta fascinante, y algún lector atento quizá ya la haya intuido.

Puesto que somos seres intrínsecamente conectados, y puesto que ni siquiera el anacoreta completo existe, si formamos una unidad profundamente entrelazada, ¿cuál es el sentido de identificar uno de esos nodos como especialmente adelantado?

Dicho de un modo particularmente claro, ¿cuál es el sentido de manifestar socialmente que alguien se ha adelantado cuando sus acciones, sus sentimientos y su propio ser están completamente embebidos en un sistema extraordinariamente interconectado?

¿Puede realmente asegurarse que el aparente logro individual de alguien premiado se debe únicamente a sus propias acciones?

¿Por qué no entender, quizá, que ese aparente adelantamiento es el resultado de una trayectoria recorrida por toda la compleja red social y que, circunstancialmente, ha terminado por hacerse visible en un determinado nodo desde el que emerge algo digno de ser destacado?

No trato de sentenciar que el premio carezca de sentido, aunque personalmente pueda inclinarme hacia esa conclusión. Trato, más bien, de mostrar que existe una necesidad profunda de cuestionar los límites de nuestra individualidad.

Aceptémoslo o no, lo que parece desprenderse de esta reflexión es que el premio, entendido como reconocimiento estrictamente individual, carece de sentido completo dentro del entramado social. Si acaso, podría comprenderse como un incentivo o, quizá con mayor precisión, como un agasajo.

También la palabra \emph{agasajo} posee una etimología sugerente. Aunque su origen exacto no esté completamente resuelto, tradicionalmente se la ha relacionado con la idea de tratar a alguien con afecto, hospitalidad o como a un compañero. Esta interpretación resulta, a mi juicio, mucho más compatible con la naturaleza profundamente interconectada de la sociedad, porque dentro de su dinamismo una persona puede decidir modificar las características de sus vínculos con otras, o al menos intentarlo.

De ese modo puede llegar a sentir a otra persona como compañera, y de ahí parece natural que el concepto haya evolucionado hasta designar el acto de entregar algo, tangible o no, a quien se considera cercano.

Aceptar un agasajo es volver recíproco el hecho de que alguien nos considere camaradas.

¿Cuál es entonces el sentido de aceptar un premio?

¿Aceptar que hemos conseguido superar el sistema de conectividad social y que, de algún modo, hemos descubierto, inventado, propuesto, escrito, hablado o pintado algo independiente de todas las conexiones que nos constituyen?

Aceptar el premio es, de algún modo, aceptar que existe el anacoreta completo.

Y ninguna persona puede, al menos en el Huniverso, alejarse lo suficiente.

\end{document}
